\scanpage{11}
% L1
C) $\unsure{\nu_0|v_0}=3,\nu_1=5,\quad\Gamma_{\nu_0,\nu_1}=\mathbb Z_{/2}\mathbb Z\times\deleted{\unsure{\mathfrak S}}\mathfrak A_5$. On a
% L2
$\alpha_0=-1,\ \alpha_1\inserted{=\alpha}$ \deleted{\unsure{sol}}solution de $\alpha^2+\alpha-1=0$, donc
% L3
$\deleted{\alpha_0+\alpha_1}\quad\alpha_0+2=1,\quad\alpha_1+2\inserted{=\alpha+2}$ inversible ($N(\alpha+2)=-1-2+4=1$,
% L4
$1/(\alpha+2)=\alpha'+2=-\alpha+1$) $\alpha_0+\alpha_1=\alpha-1$ inversible
% L5
($N(\alpha-1)=-1+1+1=1$, $1/(\alpha-1)=\alpha'-1=-\alpha-2$ ] Donc

% L6
\uline{Thm} Pour tt \deleted{hom} couple $(G,\varphi)$, $G$ forme de
% L7
$GP(1)_S$, $\varphi$ hom $\Gamma=\mathbb Z_{/2}\times\mathfrak A_5\longrightarrow G$, soit
% L8
$\alpha(\varphi)$ l'inv. de $\varphi(\rho_1)$. On trouve ainsi une
% L9
équivalence entre la cat. des $(G,\varphi)$
% L10
tels que $\forall s\in S$, \deleted{$\rho_{0s}$ et $\rho_{1s}$ soient d'ordres}\inserted{$\varphi_s\mid\mathfrak A_5$ ne soit pas trivial},
% L11
\deleted{$3$ et $5$ resp et $\rho_0$ et $\rho_1$ ne commutent pas,}
% slanted
% L11a
\margin{Laisser tomber $\mathbb Z_{/2}$ dans l'énoncé [celui sur $\mathbb Z_{/2}$, ou tt \deleted{\unsure{$\leftrightarrow$|$\Leftrightarrow$}} autre \deleted{\unsure{facteur}}\deleted{\unsure{non}} facteur, est un corollaire.]}
% L12
avec la \deleted{\unsure{ens}} cat. discrète \deleted{formée des}
% L13
définie par l'ens des $\alpha\in\Gamma(S,\unsure{\mathcal O_S|\mathscr O_S|\underline{\mathcal O}_S})$ satisfaisant
% L14
l'équation
% L15
\[\alpha^2+\alpha-1=0\]

% L16
Les représentations $\varphi$ sont triviales sur $\mathbb Z_{/2}$,
% L17
donc correspondent \unsure{aussi|ainsi} aux
% L18
$\varphi^+:\mathfrak A_5\longrightarrow G$ qui sont triviales sur
% L19
aucune fibre.
% archive sheet number: 10
